% Figure from MATH 451 notes, section 24. Compile with the notes preamble.
\begin{tikzpicture}
\begin{axis}[nfig, width=0.62\textwidth, height=0.42\textwidth, clip=false,
  xmin=0, xmax=1.0, ymin=0, ymax=11,
  xtick={0,0.5,1}, ytick=\empty, xlabel={$x$}]
  \addplot[black, very thick, domain=0:0.909, samples=150] {1/(1-x)};
  \node[font=\scriptsize, anchor=south east] at (axis cs:0.905,11) {$f = \tfrac{1}{1-x}$};
  \addplot[figB!40, thick, domain=0:0.98, samples=100] {(1-x^3)/(1-x)};
  \addplot[figB!65, thick, domain=0:0.98, samples=120] {(1-x^7)/(1-x)};
  \addplot[figB, thick, domain=0:0.96, samples=140] {(1-x^13)/(1-x)};
  \node[font=\scriptsize, figB!60!black, anchor=west] at (axis cs:0.985,2.94) {$f_2$};
  \node[font=\scriptsize, figB!80!black, anchor=west] at (axis cs:0.985,6.64) {$f_6$};
  \node[font=\scriptsize, figB, anchor=west] at (axis cs:0.965,10.3) {$f_{12}$};
\end{axis}
\end{tikzpicture}
